% !TeX root = ../report.tex
% Section 7: Limitations (the \section heading lives in report.tex).
\label{sec:limitations}

Using 379 drugs as replicates against a single trained network per arm means observed differences reflect specific network instances rather than generalisable optimiser effects, leaving null results indistinguishable from training noise. This is why the experiments did not show that SAM cannot improve rule quality: each network was trained once, and SAM+WD stopped at the 250-epoch cap. While the paired design detected small weight-decay effects ($\sim$0.007), the arms confound optimiser type, weight decay, stopping criteria, and reconstruction performance. Furthermore, the generalisation gap closed primarily via loose training fits, rendering it a weak proxy. Held-out fidelity remains bounded near 0.5 for mixed-class coverage, reaches 1.0 on single-class subsets, and relies on few test rows for $\sim$31--32\% of PDERE-GA rules and $\sim$15--18\% of PDERE-NSGA rules. No equivalence test was run, so the results cannot show that the networks are equivalent, and the largest fidelity effect (0.014) is too small to matter in practice. Finally, z-scoring pathway features across all cell lines introduced mild data leakage, though its impact is uniform across network comparisons. Biological validity and generality were also not shown: known drug--target links were almost never recovered, the KEGG test was uninformative, and nothing was tested on other datasets, architectures or extractors.